\section{Introduction}\label{sec:introduction}

Chromatic fracture reconstructs a spectrum from its local pieces and the
maps relating adjacent heights. For the sphere, one of the objects that
carries this gluing information is the overlap
\[
                         L_{n-1}L_{K(n)}S.
\]
Here \(S=S_p^\wedge\) is the derived \(p\)-complete sphere,
\(L_i\) denotes localization with respect to Johnson--Wilson theory
\(E(i)\), \(L_0\) is rationalization, and \(L_{K(n)}\) denotes Morava
\(K(n)\)-localization. Classical chromatic splitting asks for a precise
description of this overlap by lower local spheres
\citep{Hovey93,BB19}. We study the less restrictive question of whether
those spheres can produce the overlap by any finite construction.

For objects \(X_1,\ldots,X_r\) in the category of \(E(n-1)\)-local
\(S\)-modules, let
\(\Thick\{X_1,\ldots,X_r\}\) be the smallest full subcategory containing
them and closed under equivalences, finite sums, integer suspensions and
desuspensions, homotopy cofibers, and retracts
\citep[Definitions~1.2, 1.3(c), and 1.5(c)]{HS99}. Every map is \(S\)-linear.
This is the \emph{ordinary} thick subcategory: tensoring with an arbitrary
module, taking an infinite coproduct, or taking an unrestricted homotopy
limit is not an additional closure operation. The finite-assembly question
asks whether
\begin{equation}\label{eq:finite-assembly-question}
 L_{n-1}L_{K(n)}S\in
       \Thick\{L_0S,L_1S,\ldots,L_{n-1}S\}
\end{equation}
for every prime \(p\) and every \(n\geq1\). The number of cofibers and their
attaching maps may depend on \(p\) and \(n\). In particular, this question
allows nonsplit extensions and Moore corrections: a cofiber of multiplication
by a power of \(p\) is one of the permitted constructions.

\begin{theorem}\label{thm:main}
At \(p=n=3\), with \(S=S_3^\wedge\),
\[
             L_2L_{K(3)}S
                  \notin\Thick\{L_0S,L_1S,L_2S\}
\]
in the category of \(E(2)\)-local \(S\)-modules.
\end{theorem}

This disproves the ordinary finite-assembly statement
\eqref{eq:finite-assembly-question}. Its scope is the specified prime and
height. The obstruction does not depend on a proposed list of summands or
on a bound for the number of allowed cofibers.

\subsection{Splittings, finite constructions, and the prime three}

The historical splitting assertions contain more data than ordinary thick
membership. In Hovey's account of Hopkins' conjecture, the proposed exterior
classes, factorizations, and summands are distinguished from the splitting
of the canonical cofiber sequence
\citep[Conjecture~4.2]{Hovey93}. The detection result there is conditional
on chromatic splitting \citep[Theorem~4.3]{Hovey93}. For the sphere, the
weak splitting question asks only whether the canonical unit
\[
          L_{n-1}S\longrightarrow L_{n-1}L_{K(n)}S
\]
has a retraction; the finite-input formulation is stated in
\citet[Conjecture~1.1.11]{BGH2022}. Ordinary thick membership records no
distinguished map at all. A prescribed strong decomposition with its unit
would give both a finite construction and a weak splitting, but a finite
construction may have nonsplit attachments and need not split that unit.

At height two, the positive odd-prime results developed from
Shimomura--Yabe's calculation for \(p\geq5\)
\citep[Theorem~2.4]{ShimomuraYabe95}, revisited and corrected by
\citet[Theorem~7.7 and Remark~7.8]{Behrens12}.
At the prime three, Goerss--Henn--Mahowald established the prescribed
height-two splitting \citep[Theorem~5.11]{GHM14}.
Thus prime three admits a height-two splitting, while
Theorem~\ref{thm:main} obstructs even finite assembly at the next height.
The prime-two case explains why failure of a prescribed list need
not obstruct finite assembly. Beaudry disproved the original strong
wedge formula at \(n=p=2\)
\citep[Theorem~1.4]{Beaudry2017}. Beaudry--Goerss--Henn then described the
additional Moore-spectrum terms while retaining the canonical unit
retraction \citep[Theorem~1.1.6]{BGH2022}. Those extra terms are allowed by
the ordinary thick closure above. Thus the known failure of a prescribed
wedge is not by itself an obstruction to finite assembly. Conversely,
Theorem~\ref{thm:main} addresses the finite-construction question, not the
separate assertion that a particular canonical unit fails to retract.
The revised strong summand patterns and their known low-height cases are
discussed in \citet[Section~6.1]{BB19}.

The companion article \emph{Filtered chromatic splitting at generic primes}
constructs an actual filtration by lower local spheres for \(n\geq1\) and
\(p>n+1\) \citep[Theorem~1.1]{CompanionFiltered}. Its prime range excludes
the example studied here. The companion rational obstruction at height three
concerns a prescribed strong wedge for \(p\geq5\)
\citep[Corollary~1.2]{CompanionRational}; such an obstruction is compatible
with a finite construction by nonsplit cofibers. Neither companion result
supplies a step in the prime-three proof below.

The connected-marking method has earlier coheight-one precedents.
Torii's Laurent-field coefficient tower builds on Gross's monodromy
calculation and carries commuting stabilizer actions from the two heights
\citep[Theorem~3.5(1a)]{Gross79}
\citep[Theorems~2.7 and 2.9 and Corollary~2.8]{Torii03}.
Torii later constructed a localized Adams spectral sequence and detected
the degree-minus-one reduced-norm class in actual iterated localization
\citep[Theorems~4.7 and 8.1]{Torii11}. The present obstruction requires
further comparisons that retain the exceptional norm character through
ordinary finite-stage coefficients and the full lower-height deformation.

The geometric starting point is the coheight-one work of
\citet{BMR26}. For odd \(p\) and \(h\geq2\), their theorem has an
exceptional branch when \(p-1\) divides \(h-1\)
\citep[Theorem~A]{BMR26}. The calculation of the completed orientation line
assumes \(k\supseteq\mathbb F_{p^{h(h-1)}}\). At \(p=h=3\), with our choice
\(k\supseteq\mathbb F_{3^6}\), the tame determinant invariants retain a second
cohomology line with a nontrivial lower-height
norm character \citep[Theorem~D, Theorems~3.6.10 and~5.2.10, and Lemma~5.2.5]{BMR26}.
Their completed proxy does not by itself identify the
ordinary overlap used here; in particular, an invertible cofiber for the
proxy does not decide whether its attachment splits
\citep[Theorem~A and Remark~5.2.12]{BMR26}.
The shifted orientation line has a representation-theoretic antecedent in
Heyer's derived Steinberg coinvariant calculation
\citep[Corollary~5.3.4]{Heyer23}
\citep[Lemma~4.1.6 and Corollary~4.3.10]{HeyerGeo24}, used geometrically in
\citet[Sections~3.5--3.6]{BMR26}. Section~\ref{sec:analytic} computes the
required distribution orientation with the coefficient topology used here.
The proof below must carry the character from finite marking coefficients
to ordinary homotopy and retain the whole height-two deformation along the
way. These are the two places where a closed-fiber geometric calculation
alone is insufficient.

\subsection{A detector that is stable under finite construction}

The final obstruction can be stated before its coefficient calculation.
Choose a finite field \(k\) containing \(\mathbb F_{3^6}\), enlarged by a
finite extension when needed to fix the constant identifications, and let
\(E_2=E_2(k)\) be height-two Morava \(E\)-theory. Its degree-zero ring is
\(W(k)[[x]]\); the variable \(x\) is the height-two deformation parameter.
Write \(\omega=\pi_2E_2\) for the periodicity line and
\[
       Q_x=k((x)),\qquad
       \overline\omega=(\omega/3)\otimes_{k[[x]]}Q_x.
\]
An unramified quadratic field \(F/\mathbb Q_3\) lies in the height-two
endomorphism algebra. Its unit group \(G_F=\mathcal O_F^\times\) acts on
\(Q_x\) and on \(\overline\omega\). The action is semilinear: it acts on
scalars as well as on vectors.

For a spectrum \(X\), apply the exact test
\[
 \mathcal T(X)=L_{K(2)}(E_2\wedge X)/3,
 \qquad
 \mathcal V_d(X)=\pi_d\mathcal T(X)\otimes_{k[[x]]}Q_x.
\]
The smash product here is the ordinary smash product of spectra.
After \(K(2)\)-localization the three proposed generators become
\(0,0,L_{K(2)}S\), and the test of the last one has one periodicity line
in each even degree. Exactness now gives a stringent necessary condition:
if \(X\) comes from the local unit by the permitted finite operations, then
every \(\mathcal V_d(X)\) is finite-dimensional and every simple
constituent is a power of \(\overline\omega\).
Proposition~\ref{prop:constituent-test} proves this by placing those
representations in a Serre subcategory. Finite dimensionality is a
consequence on the thick closure, not an assumption about the overlap.

The coefficient calculation produces a different line,
\begin{equation}\label{eq:intro-norm-line}
 Q_x(\delta),\qquad
 \delta(g)=\overline{N_{F/\mathbb Q_3}(g)}
       \in\mathbb F_3^\times,\quad g\in G_F.
\end{equation}
The notation means that \(Q_x\) has its usual semilinear action and a chosen
basis is multiplied by the constant \(\delta(g)\). The character is
nontrivial, since the residue norm
\(\mathbb F_9^\times\to\mathbb F_3^\times\) is surjective. What matters is
not merely this residue character but its action over the field \(Q_x\).
The full group \(G_F\), including its principal units, has infinite image
on \(Q_x\) and fixed field \(k\). Together with the first Hasse invariant,
this full action excludes the norm line from every periodicity power;
Lemma~\ref{lem:generic-character} gives the argument.

\subsection{From finite markings to actual homotopy}

The coefficient calculation and its passage to homotopy solve different
problems. The first produces the norm line in ordinary continuous
cohomology. The second retains that line through the entire height-two
deformation and the filtration of actual homotopy. We describe the
interfaces between these arguments before constructing their coefficients.

\paragraph{Integral markings.}
On the height-two stratum of the height-three deformation, the algebraic
\(p\)-divisible group has a connected part of height two and an
\'etale quotient of height one. Its finite torsion kernels are formed
before passing to the stratum, so both parts are retained. Over the
one-parameter field \(K=k((t))\), let \(L^U\) be the finite \'etale
algebra of connected markings modulo an open subgroup \(U\) of the
height-two stabilizer \(J\). Put
\[
 L=\colim_{U\subset J}L^U,\qquad R=L^{\mathrm{perf}}.
\]
The groups \(P\) and \(J\) are the unextended height-three and height-two
stabilizers, respectively; both fix \(k\). The group \(P\) transports the
deformation, commuting with the marking action of \(J\). Section~\ref{sec:towers}
constructs a determinant-normalized integral marking of the \'etale
quotient. The normalization determines the norm action on the coefficients.
Relative section-sheaf full faithfulness
\citep[Corollary~3.11]{ALB25} makes this a construction of families of
markings, including their frame changes.

\paragraph{Ordinary cohomology.}
The geometry first computes cohomology with the completed perfection of
each finite stage \(L^U\). Returning to ordinary coefficients is a
separate argument: a cochain with values in \(L\) or \(R\) must have its
image in one complete finite marking or root stage. The central step in
Section~\ref{sec:decompletion} proves finite cohomology at such stages.
It uses an invariant differential to normalize Cartier's operator, whose
residue adjoint is Frobenius
\citep[Chapter~II, Section~6]{Cartier58}
\citep[Proposition~9]{Serre58Topology}. This finiteness controls primitives
of small cocycles and allows passage from root stages to completed
perfections, including continuous profinite families. A distribution
operator then removes the roots after tame sign averaging. Descent along
the remaining norm quotient gives
Theorem~\ref{thm:coefficient-calculation}: the ordinary \(P\)-cohomology
of \(L\) has four lines, in degrees \(0,1,3,4\). The upper two carry the
character \(\delta\) on \(G_F\subset J\).

\paragraph{Finite descent in spectra.}
To place this calculation in homotopy, first extend constants by a finite
\'etale \(S\)-algebra \(S_k\), whose underlying \(S\)-module is a finite
sum of copies of \(S\), and put \(Z_k=L_{K(3)}S_k\). Morava descent
expresses \(Z_k\) as the totalization of a cosimplicial spectrum.
Descendability \citep{Mathew16} bounds the composites in its error tower;
exactness of \(\mathcal T\) preserves this bound. Thus the tested descent
spectral sequence has a finite filtration on its actual abutment
\(\pi_*\mathcal T(Z_k)\), as proved in
Proposition~\ref{prop:descent-test}. On mod-\(3\) homotopy, a diagonal
summand aligns the constants from \(S_k\) and \(E_2\). Descendability
uses a thick tensor ideal; the finite-assembly detector uses the ordinary
thick closure of Section~\ref{sec:detector}.

\paragraph{The whole height-two deformation.}
The homotopy cochain terms of that spectral sequence must still be compared with
the ordinary coefficient calculation. Section~\ref{sec:homotopy} starts
with ordinary, unlocalized cooperations and expresses the completed terms
as inverse limits of their quotients by \(x^m\). Over the perfect ring
\(R\), the marked connected--\'etale extension splits. Square-zero
deformation torsors and the universal Kodaira--Spencer isomorphism
\citep[Section~5, especially Theorem~5.1 and the paragraph following
Equation~(5.2)]{Lau10} lift this splitting through
every ring \(R[x]/(x^m)\). Theorem~\ref{thm:jet-comparison} recovers the
formal-law isomorphism from the full torsion system and uses it to
evaluate ordinary cooperations. Finite Taylor expansion keeps each
evaluation in one finite coefficient stage. The resulting compatible
cochain maps retain the realized \(J\)-action through every jet and hence
on completed homotopy.

\paragraph{The surviving constituent.}
After inverting \(x\), the four rows are periodicity lines in degrees
\(0,1\) and their norm twists in degrees \(3,4\). The only possible
higher differentials go from a lower line to an upper line. The full-unit
inequivalence in Lemma~\ref{lem:generic-character} makes them zero.
Together with the finite abutment filtration, this computes the generic
diagonal test in every degree
(Proposition~\ref{prop:generic-homotopy}): it has dimension two, with one periodicity
constituent and one norm-twisted constituent. In particular, the term of
bidegree \((s,t)=(3,0)\) gives \(Q_x(\delta)\) in total degree \(-3\).
Since \(S_k\) is finite free over \(S\), finite assembly of the original
overlap would impose the constituent condition on \(\mathcal V_{-3}(Z_k)\)
as well. The norm line violates that condition.

Figure~\ref{fig:proof-chain} separates the two coefficient comparisons
from the descent and representation-theoretic inputs that let their
output survive in actual homotopy.

\begin{figure}[!htbp]
\centering
\begin{tikzpicture}[
  box/.style={draw=proofblue,rounded corners=2pt,fill=proofgray,
    align=center,text width=3.5cm,minimum height=1.0cm,
    inner sep=5pt,font=\small},
  arr/.style={-{Latex[length=2mm]},draw=proofblue,thick}
]
\node[box] (tower) at (0,0)
  {Integral marking\\towers\\Section~\ref{sec:towers}};
\node[box] (complete) at (-3,-1.65)
  {Completed\\cohomology\\Theorem~\ref{thm:completed-calculation}};
\node[box] (deform) at (3,-1.65)
  {Split deformation\\Lemma~\ref{lem:split-jet-lift}};
\node[box] (ordinary) at (-3,-3.3)
  {Ordinary cochains\\Theorem~\ref{thm:coefficient-calculation}};
\node[box] (jets) at (3,-3.3)
  {Comparison at all jets\\Theorem~\ref{thm:jet-comparison}};
\node[box] (page) at (0,-5.1)
  {Four generic rows\\Equation~\eqref{eq:obstruction-page}};
\node[box] (filtration) at (-4.7,-5.9)
  {Finite abutment\\filtration\\Proposition~\ref{prop:descent-test}};
\node[box] (units) at (4.7,-5.9)
  {Full-unit\\character exclusion\\Lemma~\ref{lem:generic-character}};
\node[box] (homotopy) at (0,-7.1)
  {Actual generic\\homotopy\\Proposition~\ref{prop:generic-homotopy}};
\node[box] (detector) at (-4.7,-8.9)
  {Finite-assembly\\detector\\Proposition~\ref{prop:constituent-test}};
\node[box] (failure) at (0,-8.9)
  {Failure of\\finite assembly\\Theorem~\ref{thm:main}};
\draw[arr] (tower) -- (complete);
\draw[arr] (tower) -- (deform);
\draw[arr] (complete) -- (ordinary);
\draw[arr] (deform) -- (jets);
\draw[arr] (ordinary) -- (jets);
\draw[arr] (jets) -- (page);
\draw[arr] (page) -- (homotopy);
\draw[arr] (filtration) -- (homotopy);
\draw[arr] (units) -- (homotopy);
\draw[arr] (homotopy) -- (failure);
\draw[arr] (detector) -- (failure);
\end{tikzpicture}
\caption{The coefficient comparisons and their passage to homotopy.
The full-unit action rules out differentials between the four generic
rows; the finite abutment filtration places their constituents in actual
homotopy. The finite-assembly detector then excludes the norm constituent.
Arrows indicate inputs to the next displayed stage.}
\label{fig:proof-chain}
\end{figure}

Section~\ref{sec:detector} proves the finite-construction test, and
Sections~\ref{sec:setup}--\ref{sec:towers} prepare the stage coefficients
and markings. Section~\ref{sec:finite-resolutions} supplies the finite
analytic resolution needed in Sections~\ref{sec:analytic}--\ref{sec:decompletion}.
Section~\ref{sec:homotopy} proves convergence and identifies the ordinary
cooperations, which Section~\ref{sec:jets} compares through every jet.
Section~\ref{sec:obstruction} completes the argument.
\FloatBarrier
